
\documentclass{beamer}
\usetheme{Madrid}
\usecolortheme{seahorse}
\usepackage{amsmath, amsfonts}
\usepackage{graphicx}
\usepackage{tikz}
\usepackage{qcircuit}
\usepackage{caption}

\title[QML with NN]{Quantum Machine Learning with Neural Networks}
\author{Morten Hjorth-Jensen}
\institute{University of Oslo}
\date{\today}

\begin{document}

\begin{frame}
  \titlepage
\end{frame}

\begin{frame}{Outline}
  \tableofcontents
\end{frame}

\section{Overview of QML}
\begin{frame}{Quantum Machine Learning Landscape}
  \begin{itemize}
    \item Quantum Data and Classical Data: \( |\psi\rangle \) vs vectors \( \mathbf{x} \in \mathbb{R}^n \)
    \item Categories:
    \begin{itemize}
        \item Quantum-enhanced classical ML
        \item Quantum-native ML algorithms
        \item Hybrid quantum-classical ML
    \end{itemize}
    \item Resource-aware quantum learning models
  \end{itemize}
\end{frame}

\section{Quantum Computing Recap}
\begin{frame}{Quantum States and Gates}
  \begin{itemize}
    \item Qubit: \( |\psi\rangle = \alpha|0\rangle + \beta|1\rangle \), \( |\alpha|^2 + |\beta|^2 = 1 \)
    \item Common Gates:
    \[
        H = \frac{1}{\sqrt{2}}\begin{bmatrix}1 & 1\\ 1 & -1\end{bmatrix}, \quad
        X = \begin{bmatrix}0 & 1\\ 1 & 0\end{bmatrix}
    \]
    \item Entanglement via CNOT and multi-qubit systems
  \end{itemize}
\end{frame}

\begin{frame}{Quantum Circuit Example}
\[\Qcircuit @C=1em @R=.7em {
  \lstick{|0\rangle} & \gate{H} & \ctrl{1} & \qw \\
  \lstick{|0\rangle} & \qw      & \targ   & \qw
}\]
\begin{itemize}
  \item Bell state: \( \frac{1}{\sqrt{2}}(|00\rangle + |11\rangle) \)
\end{itemize}
\end{frame}

\section{Neural Networks and Representational Capacity}
\begin{frame}{Classical Neural Networks Review}
  \begin{itemize}
    \item Universal approximation theorem
    \item Layer-wise structure: \( y = \sigma(Wx + b) \)
    \item Deep networks and expressivity
    \item Limitations in high-dimensional feature space
  \end{itemize}
\end{frame}

\section{Quantum Neural Networks (QNN)}
\begin{frame}{What is a QNN?}
  \begin{itemize}
    \item Quantum analog of classical NNs using PQCs
    \item Data encoding + entangling + variational layer
    \item Output via observable measurements
  \end{itemize}
\end{frame}

\begin{frame}{Generic QNN Architecture}
\begin{center}
%  \includegraphics[width=0.85\linewidth]{qnn_diagram.png}
\end{center}
\captionof{figure}{QNN: Feature map + variational layers + measurement}
\end{frame}

\section{Variational Quantum Circuits (VQCs)}
\begin{frame}{Variational Quantum Circuits}
  \begin{itemize}
    \item PQCs with parameters \( \theta \)
    \item Optimize: \( C(\theta) = \langle \psi(\theta)| \hat{H} |\psi(\theta)\rangle \)
  \end{itemize}
\end{frame}

\begin{frame}{Circuit Example: Ansatz Layer}
\[\Qcircuit @C=1em @R=.7em {
  \lstick{|x_1\rangle} & \gate{R_Y(x_1)} & \multigate{1}{Entangle} & \gate{R_Y(\theta_1)} & \qw \\
  \lstick{|x_2\rangle} & \gate{R_Y(x_2)} & \ghost{Entangle}       & \gate{R_Y(\theta_2)} & \qw
}\]
\end{frame}

\section{Training QNNs}
\begin{frame}{Training and Optimization}
  \begin{itemize}
    \item Gradient via Parameter Shift Rule:
    \[ \frac{\partial \langle O \rangle}{\partial \theta} =
    \frac{\langle O \rangle_{\theta + \pi/2} - \langle O \rangle_{\theta - \pi/2}}{2} \]
  \end{itemize}
\end{frame}

\section{Applications}
\begin{frame}{Applications of QML + NN}
  \begin{itemize}
    \item Quantum-enhanced classification
    \item QGANs, Quantum autoencoders
    \item Anomaly detection, Quantum RL
  \end{itemize}
\end{frame}

\section{Open Challenges}
\begin{frame}{Challenges in QNNs}
  \begin{itemize}
    \item Hardware limitations, barren plateaus
    \item Hybrid training instability
  \end{itemize}
\end{frame}

\section{Conclusion and Future Work}
\begin{frame}{Future Directions}
  \begin{itemize}
    \item Efficient hybrid architectures
    \item Adaptive QNNs and better ansätze
  \end{itemize}
\end{frame}

\begin{frame}{Thank You!}
  \centering
  \Huge Questions?
\end{frame}


% Section: Quantum Convolutional Neural Networks
\section{Quantum Convolutional Neural Networks (QCNNs)}
\begin{frame}{QCNN Architecture}
  \begin{itemize}
    \item Inspired by classical CNNs for hierarchical feature extraction.
    \item Layers include quantum convolution and pooling operations.
    \item Reduce qubit counts while preserving quantum information.
  \end{itemize}
\end{frame}

\begin{frame}{QCNN Example Circuit}
\begin{itemize}
  \item Quantum convolution applies parameterized gates to pairs of qubits.
  \item Pooling reduces the system size, e.g., via measurement or entanglement filtering.
\end{itemize}
\begin{center}
%  \includegraphics[width=0.75\linewidth]{qcnn_diagram.png}
\end{center}
\captionof{figure}{Illustrative QCNN circuit with convolution and pooling layers.}
\end{frame}

% Section: Quantum Generative Adversarial Networks
\section{Quantum GANs (QGANs)}
\begin{frame}{Quantum GAN Framework}
  \begin{itemize}
    \item Combines quantum generator and classical or quantum discriminator.
    \item Generator learns to produce quantum states matching target distribution.
    \item Applications: quantum state preparation, data augmentation.
  \end{itemize}
\end{frame}

\begin{frame}{QGAN Architecture}
\begin{center}
%  \includegraphics[width=0.8\linewidth]{qgan_diagram.png}
\end{center}
\captionof{figure}{Example QGAN setup with a quantum generator and classical discriminator.}
\end{frame}

% Section: Benchmark Results
\section{Benchmark Results and Performance}
\begin{frame}{Benchmarks in QML Research}
  \begin{itemize}
    \item Performance metrics: classification accuracy, fidelity, loss convergence.
    \item Datasets: quantum-enhanced MNIST, quantum chemistry datasets.
    \item Notable results:
    \begin{itemize}
      \item QNN classifiers outperform classical on small quantum datasets.
      \item QCNNs demonstrate improved generalization with fewer qubits.
      \item QGANs achieve state fidelity > 95\% on target distributions.
    \end{itemize}
  \end{itemize}
\end{frame}

\begin{frame}{Example Benchmark Table}
\begin{center}
\begin{tabular}{|c|c|c|c|}
\hline
Model & Dataset & Accuracy (\%) & Qubits Used \\
\hline
QNN & Iris (Quantum Enc.) & 94.3 & 4 \\
QCNN & Quantum MNIST & 89.7 & 6 \\
QGAN & Quantum States & 95.5 (Fidelity) & 5 \\
\hline
\end{tabular}
\end{center}
\captionof{table}{Sample benchmark results for QNN, QCNN, and QGAN models.}
\end{frame}

\end{document}
